% ===================================================================
%  সারণি নং 8 — ফসল। — শিকার। — আঙুরের ফসল। — মাছ ধরা।
%
%  Ceci n'est PAS une transcription : c'est une traduction, faite en
%  2026, en bengali d'aujourd'hui. Voir texto/bn/10-sarani-01.tex
%  pour la consigne, qui vaut pour les seize fichiers.
% ===================================================================

%%K t08-noto-1 noto
\VUnotes{143.2mm}{%
\textit{(*) কারণ এই শব্দটি অস্পষ্ট: শণের ফসল? আঙুরের? আপেলের? হয়তো ভালো হত
«} \VUgras{mesono} \textit{» কাজে লাগানো, যা} Mondo \textit{XIV, পৃ. 242-এ
প্রস্তাবিত। কিন্তু এ পর্যন্ত সেই নতুন ধাতু আকাদেমি গ্রহণ করেনি এবং তা ইদো
আন্দোলনের সাধারণ নিয়ম অনুসারে আলোচনার অপেক্ষায় আছে।}%
}

%%K t08-apar-1 apar
\VUsaut{5.0mm}
\VUtitre{15.0pt}{60}{সারণি নং 8}
\VUsaut{3.0mm}
\VUcentre{13.2pt}{90}{\textit{প্রথম দৃশ্য।}}
\VUsaut{3.0mm}
\VUpk{10.2pt}{40}{ফসল (*)।}
\VUsaut{4.0mm}
\VUpk{10.2pt}{40}{মাঠের রূপ।}
\VUsaut{3.0mm}

%%K t08-c1-01-1 p
1. --- \VUgras{গ্রীষ্মের} সঙ্গে মাঠ আরও একবার নিজের রূপ বদলে ফেলেছে। সিতায়
সঁপে দেওয়া বীজ ফুটে উঠেছে; মাটি থেকে একটি ছোট্ট সবুজ চারা বেরিয়ে নিজের শিষ
দিয়েছে। দানা বেঁধেছে, তারপর পেকেছে, আর এখন ফসল কাটা ছাড়া কিছু বাকি নেই।

%%K t08-c1-02-1 p
2. --- \VUgras{বুনো ফুল} ফুটেছে। \VUgras{কর্নফ্লাওয়ার} \textsuperscript{(1)},
\VUgras{পপি} \textsuperscript{(2)} ও \VUgras{ডিজিটালিস} \textsuperscript{(3)}
গমখেতের একরঙা পটে নিজেদের টাটকা \VUgras{ফুলের} প্রসন্ন বৈপরীত্য রাখে।

%%K t08-c1-03-1 p
3. --- ফলগাছে পাতা ধরেছে; ফল পেকেছে। ছোট্ট \VUgras{চেরিগাছ}
\textsuperscript{(4)} সুন্দর \VUgras{চেরিতে} \textsuperscript{(5)} লদপদ, যা
ছেলেমেয়েরা পেড়ে সাগ্রহে খায়।

%%K t08-c1-04-1 p
4. --- পাল এখন সারা দিন বাইরে কাটাতে পারে।

%%K t08-c1-04-2 p
\VUgras{মেষশাবকেরা} \textsuperscript{(6)} বেড়ে সুন্দর \VUgras{ভেড়া}
\textsuperscript{(7)} ও সুন্দর \VUgras{মেষ} \textsuperscript{(8)} হয়েছে, ঘন
পশমওয়ালা। তারা \VUgras{চারণভূমির} \textsuperscript{(9)} \VUgras{ঘাস} শান্তিতে
চরে, যার উপরে \VUgras{ডেইজি} ছড়ানো।

%%K t08-c1-04-3 p
শিগগিরই এই জন্তুদের \VUgras{পশম} ছাঁটা হবে, যা থেকে আমাদের জামাকাপড়ের কাপড়
তৈরি হয়।

%%K t08-tit-2 sub
\VUsaut{5.0mm}
\VUpk{10.2pt}{40}{রাখাল।}
\VUsaut{3.4mm}

%%K t08-c1-05-1 p
5. --- \VUgras{রাখাল} \textsuperscript{(10)} তাদের দিকে বেশি নজর দিচ্ছে বলে মনে
হয় না। তার একটিই চিন্তা, দিগন্তে সরে যাওয়া ঝড়, আর \VUgras{মেষপাল}
\textsuperscript{(13)}, যে নিজের \VUgras{সুরাহি} \textsuperscript{(14)} ঠোঁটে
তুলেছে, তার চেয়েও বেশি নির্ভাবনা।

%%K t08-c1-06-1 p
6. --- গরমটা দমবন্ধ করা; কারণ \VUgras{কুকুরও} \textsuperscript{(15)} ঠান্ডা হতে
আসে; সে \VUgras{নালার} \textsuperscript{(16)} টাটকা জল চাটে ও একটি বেচারা ছোট্ট
\VUgras{ব্যাঙকে} \textsuperscript{(17)} চমকে দেয়, যে তীরের ঘাসে গুটিয়ে ছিল। সে
সেই স্বচ্ছ জলে ঝাঁপ দেয় যেখানে \VUgras{শালুক} \textsuperscript{(18)} ফুটেছে।

%%K t08-c1-07-1 p
7. --- একটি \VUgras{খঞ্জন} \textsuperscript{(19)} সেই ছোট্ট \VUgras{ঝোরার}
\textsuperscript{(20)} পাশে নাইছে যা দুই পাথরের ফাঁক থেকে ফোটে ও
\VUgras{কাঁটাঝোপ} \textsuperscript{(21)} এবং \VUgras{কুলগাছের ঝোপের}
\textsuperscript{(22)} নিচে লুকিয়ে পড়ে।

%%K t08-tit-3 sub
\VUsaut{5.0mm}
\VUpk{10.2pt}{40}{ফসল কাটা।}
\VUsaut{3.4mm}

%%K t08-c1-08-1 p
8. --- মাঠের ছেলেমেয়েদের কাছে গম কাটা বড় আনন্দের সময়। তারা প্রসন্ন
\VUgras{ডিগবাজি} \textsuperscript{(24)} খায় \VUgras{খড়ের স্তূপে}
\textsuperscript{(23)}, \VUgras{কাটিয়েদের} \textsuperscript{(26)} সাহায্য করে,
আর যখন তারা \VUgras{গমখেত} \textsuperscript{(27)} কাটে নিজেদের \VUgras{কাস্তে}
\textsuperscript{(25)} দিয়ে, যা \VUgras{শানপাথরে} \textsuperscript{(30)} বেশ ধার
দেওয়া, তখন ছেলেমেয়েরা গম জড়ো করে \VUgras{আঁটি} \textsuperscript{(31)} বা
\VUgras{ছোট বোঝা} \textsuperscript{(32)} বাঁধে।

%%K t08-c1-09-1 p
9. --- কখনও কখনও তাদের মধ্যে বড়দের \VUgras{কাঁচি-কাস্তে}
\textsuperscript{(12)} দিয়ে সেই \VUgras{সোনালি ডাঁটা} \textsuperscript{(28)}
কাটতে দেওয়া হয় যা দানায় ভরা ভারী \VUgras{শিষ} \textsuperscript{(29)} তুলে
আছে; আর ছোটরা, \VUgras{গবাদির} \textsuperscript{(33)} রক্ষণাবেক্ষণ করতে করতে, যা
\VUgras{মাঠে} \textsuperscript{(34)} বড় \VUgras{লিন্ডেন গাছের}
\textsuperscript{(35)} ছায়ায় চরে, নিজেদের \VUgras{জালে} \textsuperscript{(36)}
সুন্দর \VUgras{প্রজাপতি} ধরে।

%%K t08-c1-09-2 p
কেউ কেউ তো \VUgras{গাড়িতে} \textsuperscript{(37)} উঠে সেই আঁটিগুলি সাজায় যা
তাদের \VUgras{আঁকশিতে} \textsuperscript{(38)} তুলে দেওয়া হয়।

%%K t08-c1-10-1 p
10. --- সারা \VUgras{সমতলে} \textsuperscript{(39)}, যেখান থেকে
\VUgras{ভরতপাখিরা} \textsuperscript{(41)} উড়ে যায়, বড় সাড়া। সবাই ফসলে লেগে।
কিন্তু যেখানে গরিবেরা পুরনো যন্ত্রপাতিই কাজে লাগায় --- \VUgras{কাস্তে,
কাঁচি-কাস্তে} ও \VUgras{মুগুর} \textsuperscript{(40)} ---, সেখানে ধনী ভূস্বামীরা
নিজেদের গম বা নিজেদের \VUgras{রাই} \textsuperscript{(43)} \VUgras{কাটার
যন্ত্র} \textsuperscript{(42)} দিয়ে কাটান। তারপর, আঁটিগুলি খামারঘরে না বয়েই,
সেখানেই বড় \VUgras{গাদা} \textsuperscript{(44)} দেওয়া হয়।

%%K t08-c1-11-1 p
11. --- তারপর একটি \VUgras{মাড়াইয়ের যন্ত্র} \textsuperscript{(47)} আনা হয়,
যাকে একটি \VUgras{চলন্ত ইঞ্জিন} \textsuperscript{(45)} একটি \VUgras{বেল্টের}
\textsuperscript{(46)} মাধ্যমে চালায়, আর এইভাবে দানা \VUgras{খড়} থেকে আলাদা হয়ে
যায়, সেই খেতেই যেখানে তা বোনা হয়েছিল।

%%K t08-tit-4 sub
\VUsaut{5.0mm}
\VUpk{10.2pt}{40}{ঝড়।}
\VUsaut{3.4mm}

%%K t08-c1-12-1 p
12. --- ফসলের সময় প্রায়ই প্রচণ্ড \VUgras{ঝড়} \textsuperscript{(53)} ভেঙে পড়ে।
প্রথমে দূর থেকে \VUgras{মেঘের গর্জন} শোনা যায়; একটি প্রবল \VUgras{পশ্চিমা
বাতাস} \textsuperscript{(54)} ওঠে ও \VUgras{বনের} \textsuperscript{(49)} সবচেয়ে
বড় গাছগুলিকে গোঙাতে গোঙাতে নুইয়ে দেয়। কালো \VUgras{মেঘ}
\textsuperscript{(56)} আকাশে চড়ে আসে; তাদের চিরে দেয় \VUgras{বিদ্যুতের}
\textsuperscript{(55)} ধাঁধানো আঁকাবাঁকা রেখা। \VUgras{বৃষ্টি} আর কখনও
\VUgras{শিলা} পড়তে শুরু করে, আর কাটার জন্য তৈরি ফসল শুইয়ে দেয়।

%%K t08-c1-13-1 p
13. --- প্রায়ই বাজ কোনো ইমারতে আগুন ধরিয়ে দেয়। গত বছর, উদাহরণস্বরূপ,
\VUgras{বায়ুচক্কির} \textsuperscript{(50)} ছাদ পুড়ে ছাই হয়ে গিয়েছিল আর তার
দুটি \VUgras{ডানা} \textsuperscript{(51)} চুরমার হয়ে গিয়েছিল।

%%K t08-c1-14-1 p
14. --- কয়েক ঘণ্টা পরে শান্তি ফিরে আসে। দিগন্তে একটি ঝলমলে \VUgras{রামধনু}
\textsuperscript{(52)} দেখা দেয়, আর প্রকৃতি সেই জীবন আবার তুলে নেয় যা কিছুক্ষণের
জন্য থেমে গিয়েছিল। যে গ্রামবাসীরা \VUgras{কুঁড়েঘরে} \textsuperscript{(48)}
এসে লুকিয়েছিল, তারা কাজে ফেরে ও এইমাত্র হওয়া ক্ষতি সারাতে সাহসে লেগে
পড়ে।

%%K t08-tit-5 sub
\VUsaut{6.0mm}
\VUcentre{13.2pt}{90}{\textit{দ্বিতীয় দৃশ্য।}}
\VUsaut{3.6mm}

%%K t08-tit-6 sub
\VUpk{10.2pt}{40}{শিকার। --- আঙুরের ফসল।}
\VUsaut{1.4mm}
\VUpk{10.2pt}{40}{মাছ ধরা।}
\VUsaut{4.0mm}

%%K t08-c2-01-1 p
1. --- \VUgras{আগস্টের} শেষে বা \VUgras{সেপ্টেম্বরের} মাঝামাঝি, অঞ্চল অনুসারে,
শিকারের মরশুম খোলে। সূর্যের প্রথম কিরণ \VUgras{টিকটিকিদের}
\textsuperscript{(2)} গর্ত থেকে বের করে আনতে না আনতেই \VUgras{শিকারি}
\textsuperscript{(5)} নিজের \VUgras{কুকুর} \textsuperscript{(4)} নিয়ে বেরিয়ে
পড়ে।

%%K t08-c2-02-1 p
2. --- সে মোটা কাপড়ের নিজের মজবুত \VUgras{শিকারি পোশাক} \textsuperscript{(9)}
পরেছে, আর চামড়ার \VUgras{পটি}, যা তাকে সেই ভেজা ঘাসে হাঁটতে দেবে যেখানে
\VUgras{কোলাব্যাঙ} \textsuperscript{(1)} লুকিয়ে থাকে; সে নিজের \VUgras{বন্দুক}
\textsuperscript{(6)} ও নিজের \VUgras{শিকার-থলি} \textsuperscript{(10)} নিয়েছে,
আর চলে গেছে।

%%K t08-c2-03-1 p
3. --- শিকারের উত্তেজনা তাকে একটি আঙুরবাগানের মাঝে এনে ফেলে যেখানে আঙুরের ফসল
পুরো জোরে চলছে। আঙুরচাষির ছেলেমেয়েরা, যারা সাধারণত রাস্তার ধারে ছাগল চরায়,
আজ তাদের যেখানে ইচ্ছা চরতে দিচ্ছে, আর একটি বুড়ো \VUgras{পাঁঠাকে}
\textsuperscript{(3)} খোঁটায় বেঁধে --- যে নিজের স্বাধীনতার সুযোগ নিয়ে নিশ্চয়ই
পালাত --- তারাও নিজেদের ভাগের মজা নিচ্ছে। একজন একটি \VUgras{আঙুরের ঝুড়ি}
\textsuperscript{(12)} থেকে একটি থোকা নেয় আর তার \VUgras{রস} \textsuperscript{(14)}
একটি \VUgras{পেয়ালায়} \textsuperscript{(13)} নিংড়ে মস্ত (গাঁজন ছাড়া মদ)
বানাতে মন বসায়। অন্যজন নিজের মাথায় \VUgras{পাতার মুকুট} \textsuperscript{(15)}
পরছে যা সে এইমাত্র গেঁথেছে। তাদের পাশে ধৃষ্ট \VUgras{চড়ুইরা}
\textsuperscript{(16)} আঙুর চুরি করতে ঘানি পর্যন্ত চলে আসে।

%%K t08-c2-04-1 p
4. --- ছেলেমেয়েরা যতক্ষণ মজা করে, লোকেরা কাজ করে। \VUgras{আঙুর-তোলকেরা}
\textsuperscript{(18)} আঙুর \VUgras{পিঠের ঝুড়িতে} \textsuperscript{(17)}
ভাঁড়ারঘরে নিয়ে যায়; একজন \VUgras{পিপেকার} \textsuperscript{(19)}
\VUgras{পিপে} \textsuperscript{(20)} সারায়, দেখে যে \VUgras{ফলক}
\textsuperscript{(22)} ও \VUgras{ঢাকনা} \textsuperscript{(21)} ঠিক আছে, আর ভাঙা
\VUgras{বেড়} \textsuperscript{(23)} বদলায়। সেই বড় \VUgras{চৌবাচ্চাও}
\textsuperscript{(25)} সে-ই বানিয়েছিল, যাতে গাঁজন তুলতে \VUgras{আঙুর}
\textsuperscript{(26)} ঢালা হয়।

%%K t08-c2-05-1 p
5. --- গাঁজন হয়ে গেলে মদ বের করে নেওয়া হয় আর বাকিটা \VUgras{ঘানিতে}
\textsuperscript{(28)} রাখা হয়; বড় \VUgras{পেঁচকে} \textsuperscript{(29)} ধীরে
ধীরে কষে রস নিংড়ানো হয়, যা একটি \VUgras{গামলায়} \textsuperscript{(32)} বয়
আর আরও \VUgras{মদ} \textsuperscript{(31)} দেয়।

%%K t08-tit-8 sub
\VUsaut{6.0mm}
\VUpk{10.2pt}{40}{বিয়ার ও সিডার।}
\VUsaut{3.4mm}

%%K t08-c2-06-1 p
6. --- \VUgras{ভাঁড়ারঘরের} \textsuperscript{(27)} দেয়ালে \VUgras{হপসের}
\textsuperscript{(30)} একটি ডাল ওঠে। এখানে সেটি কেবল সাজানোর গাছ, কিন্তু কিছু
দেশে, যেখানে আঙুরলতা খুব কম হয়, \VUgras{বিয়ার} বানাতে হপস চাষ করা হয়।

%%K t08-c2-07-1 p
7. --- ছোট্ট \VUgras{আপেলগাছ} \textsuperscript{(33)} আপেলে লদপদ, যা পাকলে
চমৎকার \VUgras{সিডার} হবে। নিজের \VUgras{খাঁচায়} \textsuperscript{(34)}
\VUgras{ক্যানারি} \textsuperscript{(35)} ও \VUgras{গোল্ডফিঞ্চ}
\textsuperscript{(36)} সেই চড়ুইদের হিংসায় দেখে যারা স্বাধীন হয়ে লাফালাফি
করছে।

%%K t08-c2-08-1 p
8. --- যতদূর চোখ যায়, পাহাড়ের ঢালে, \VUgras{আঙুরলতার খোঁটার}
\textsuperscript{(40)} সারি দাঁড়িয়ে, যা চমৎকার \VUgras{আঙুরগাছ}
\textsuperscript{(39)} ধরে আছে, \VUgras{থোকায়} ভরা।

%%K t08-c2-08-2 p
সারা বছর \VUgras{আঙুরচাষি} \textsuperscript{(42)} নিজের \VUgras{আঙুরবাগানের}
\textsuperscript{(38)} যত্নে খেটেছে, আর এখন সে ভালো ফসলে তার পুরস্কার পাচ্ছে।
\VUgras{তোলকিনীরা} \textsuperscript{(41)} সেই চৌবাচ্চাগুলি ভরছে যা
\VUgras{খচ্চরে} \textsuperscript{(43)} টানা গাড়িতে রাখা, আর সবাই খুশি।

%%K t08-tit-9 sub
\VUsaut{5.0mm}
\VUpk{10.2pt}{40}{মাছ ধরা।}
\VUsaut{3.4mm}

%%K t08-c2-09-1 p
9. --- একই কথা \VUgras{ছিপওয়ালাদের} \textsuperscript{(44)}, যারা ভোর হতেই
নিজেদের \VUgras{ছিপ} \textsuperscript{(45)} নিয়ে জলের ধারে এসে জমেছে।

%%K t08-c2-09-2 p
\VUgras{মাছ} \textsuperscript{(47)} বঁড়শি ধরে ফেলে যেই তারা নিজেদের
\VUgras{সুতো} \textsuperscript{(46)} জলে ফেলে, আর তাদের \VUgras{টোপ} বদলানোর
সময়ও মেলে না, তার আগেই সুতোর ডগায় একটি নতুন শিকার ছটফট করতে থাকে।

%%K t08-c2-09-3 p
তবু \VUgras{জালওয়ালা জেলে} \textsuperscript{(48)}, যে নিজের \VUgras{নৌকা}
\textsuperscript{(49)} থেকে \VUgras{নদীর} \textsuperscript{(50)} মাঝখানে জাল
ফেলে, অনেক বেশি ধরে।

%%K t08-c2-10-1 p
10. --- শরৎ ভালো বেড়ানোর মরশুম: পায়ে হেঁটে, \VUgras{ঘোড়ায়}
\textsuperscript{(52)} চড়ে, \VUgras{সাইকেলে} \textsuperscript{(53)},
\VUgras{মোটরে} \textsuperscript{(54)}। \VUgras{রাস্তা} \textsuperscript{(51)}
পরিষ্কার আর লোকে শেষ ভালো দিনগুলির পুরো সদ্ব্যবহার করে, কারণ
\VUgras{আবাবিলরা} \textsuperscript{(56)} আগেই ছাদে জমতে শুরু করেছে, গরম দেশে উড়ে
যাওয়ার জন্য। বুড়ো \VUgras{আখরোট গাছের} \textsuperscript{(57)} পাতা হলুদ হয়ে
আসছে, \VUgras{আখরোট} ঝরছে, আর একটি \VUgras{ঝোপে} \textsuperscript{(58)} দাঁড়ানো
উঁচু \VUgras{গাছ}, যা \VUgras{উঁচু জমিতে} \textsuperscript{(59)}, শিগগিরই ন্যাড়া
হয়ে যাবে।

%%K t08-c2-11-1 p
11. --- \VUgras{সূর্য} \textsuperscript{(60)} দেরিতে ওঠে, দিন ছোট হয়ে আসে,
\VUgras{সকালে} বেশ ভালোই ঠান্ডা লাগতে শুরু করে, আর \VUgras{বনমোরগের}
\textsuperscript{(61)} ঝাঁক এখন থেকেই আমাদের নির্মম আবহাওয়া ছেড়ে চলে যাচ্ছে।
\VUsaut{6.0mm}
\VUfilet{22mm}
